% !TEX program = lualatex
\documentclass[rodape]{synergiakickoff}
\dxiprojeto{SYNERGIA · Integração de IA · Equipe 2}

\begin{document}

\dxicapa
  {IFAM · INOVA}
  {\includegraphics[width=13.2cm]{marca/synergia-horizontal-branca.png}}
  {IA para ampliar a análise de indicadores de Suprimentos.}

\dxiequipe[Equipe e papéis]{%
  \dximembro[fotos/card-lucas]{P.O. LG}{Lucas Costa}
  \dximembro[fotos/card-gabriel]{P.O. · Fullstack}{Gabriel Fernandes}
  \dximembro[fotos/card-rebecca]{Frontend · Design}{Rebecca Xavier}
  \dximembro[fotos/card-marcelo]{DevOps · QA}{Marcelo Uchôa}
  \dximembro[fotos/card-carlos]{RPA}{Carlos Eduardo}
  \dximembro[fotos/card-gustavo]{Backend}{Gustavo Nunes}
}

\begin{dxiconteudo}{CONTEXTO DO PROJETO}
\vspace{-0.40cm}
\dxicol{13.7}{%
  {\dximont\bfseries\fontsize{16}{21}\selectfont\color{lgred}O desafio}

  \vspace{0.35em}
  Informações de plano, produção, recebimento e qualidade chegam de fontes
  distintas. A equipe precisa cruzar Workorders, lotes e seriais para localizar
  divergências e preparar o acompanhamento operacional.

  \vspace{0.65em}
  O levantamento inicial citou \textbf{44 horas por mês}, mas esse indicador
  ainda depende de medição e validação do cliente. Nenhum ganho de FTE é
  declarado nesta proposta.
}%
\dxicol{13.1}{%
  {\dximont\bfseries\fontsize{16}{21}\selectfont\color{hanaroverde}O SYNERGIA hoje}

  \vspace{0.35em}
  \begin{dxilista}[hanaroverde]
    \item Aplicação em Angular, FastAPI e PostgreSQL.
    \item Ingestão manual de CSV e XLSX enquanto conectores não estão disponíveis.
    \item Validação, normalização, consolidação e regras OQC determinísticas.
    \item Dashboard, relatórios, auditoria e decisão humana.
  \end{dxilista}

  \vspace{0.45em}
  {\dxiCorpoP\color{dxinavy!85}A IA entra como apoio à interpretação. As regras
  e a decisão continuam sob controle do sistema e das pessoas.}
}
\end{dxiconteudo}

\begin{dxiconteudo}{DUAS OPORTUNIDADES DE IA}
\vspace{-0.35cm}
\dxicol{13.4}{%
  {\dximont\bfseries\fontsize{18}{24}\selectfont\color{lgred}1. Qualidade dos dados}

  \vspace{0.4em}
  Transformar erros, avisos e divergências já detectados em uma explicação
  objetiva, com impacto, evidências e próxima verificação sugerida.

  \vspace{0.8em}
  {\dximont\bfseries\fontsize{18}{24}\selectfont\color{lgred}2. Investigação operacional}

  \vspace{0.4em}
  Reunir o contexto de uma execução e produzir um resumo rastreável para apoiar
  a análise de pendências e Workorders.
}%
\dxicol{13.4}{%
  {\dximont\bfseries\fontsize{16}{21}\selectfont\color{hanaroverde}Princípios de integração}

  \vspace{0.35em}
  \begin{dxilista}[hanaroverde]
    \item Modelo executado localmente, sem envio de dados a serviços externos.
    \item Ferramentas de leitura com permissões e escopo organizacional existentes.
    \item Saída estruturada, com evidências e indicação de informação insuficiente.
    \item Nenhuma escrita, aprovação, notificação ou RPA comandada pelo modelo.
    \item Validação humana antes de qualquer ação posterior.
  \end{dxilista}
}
\end{dxiconteudo}

\begin{dxiconteudo}{PROPOSTA 1 · ASSISTENTE DE QUALIDADE}
\vspace{-0.50cm}
\dxicol{12.9}{%
  {\dximont\bfseries\color{lgred}D.E.E.P AI}\par
  \begin{dxilista}[dxired]
    \item \textbf{Define:} reduzir o tempo de leitura das ocorrências de qualidade.
    \item \textbf{Explore:} usar validações, divergências e proveniência autorizada.
    \item \textbf{Experiment:} comparar explicações com uma baseline determinística.
    \item \textbf{Productionize:} integrar somente após avaliação e aceite do negócio.
  \end{dxilista}
}%
\dxicol{13.9}{%
  {\dximont\bfseries\color{hanaroverde}Experiência proposta}\par
  \begin{dxilista}[hanaroverde]
    \item O operador abre uma execução com erro ou advertência.
    \item O backend monta um contexto mínimo com IDs de evidência válidos.
    \item A IA explica a ocorrência e sugere a próxima conferência.
    \item A interface exibe fonte, limitação e necessidade de revisão humana.
  \end{dxilista}

  \vspace{0.45em}
  {\dxiCorpoP\color{dxinavy!85}Valor esperado: tornar o diagnóstico mais claro
  sem refazer a classificação já produzida pelas regras do SYNERGIA.}
}
\end{dxiconteudo}

\begin{dxiconteudo}{INTEGRAÇÃO · ASSISTENTE DE QUALIDADE}
\vspace{-0.40cm}
\dxicol{13.6}{%
  {\dximont\bfseries\fontsize{16}{21}\selectfont\color{lgred}Fluxo técnico}

  \vspace{0.3em}
  \begin{dxilista}[dxired]
    \item FastAPI seleciona somente ocorrências autorizadas.
    \item Adaptador remove campos livres e dados desnecessários.
    \item Runtime local recebe contexto, prompt e schema versionados.
    \item Validador rejeita evidência inexistente, campo extra ou contradição.
    \item Angular apresenta a resposta ao lado do dado determinístico.
  \end{dxilista}
}%
\dxicol{13.2}{%
  {\dximont\bfseries\fontsize{16}{21}\selectfont\color{hanaroverde}Critérios para avançar}

  \vspace{0.3em}
  \begin{dxilista}[hanaroverde]
    \item Conjunto sintético e gabarito revisados por pessoa independente.
    \item Comparação com baseline nos mesmos casos.
    \item Evidências válidas e abstenção quando faltar contexto.
    \item Latência e memória compatíveis com o ambiente definido.
    \item Aprovação técnica, de segurança e do cliente.
  \end{dxilista}
}
\end{dxiconteudo}

\begin{dxiconteudo}{PROPOSTA 2 · ASSISTENTE DE INVESTIGAÇÃO}
\vspace{-0.50cm}
\dxicol{12.9}{%
  {\dximont\bfseries\color{lgred}D.E.E.P AI}\par
  \begin{dxilista}[dxired]
    \item \textbf{Define:} reduzir consultas dispersas durante uma investigação.
    \item \textbf{Explore:} consultar execução, pendências e Workorders autorizados.
    \item \textbf{Experiment:} avaliar resumo, evidências e seleção de ferramentas.
    \item \textbf{Productionize:} liberar por etapas, começando com contexto pronto.
  \end{dxilista}
}%
\dxicol{13.9}{%
  {\dximont\bfseries\color{hanaroverde}Experiência proposta}\par
  \begin{dxilista}[hanaroverde]
    \item O usuário escolhe uma execução dentro de sua organização.
    \item O sistema recupera fatos pelas permissões da sessão.
    \item A IA organiza achados, perguntas em aberto e próximos passos humanos.
    \item Cada afirmação aponta para uma evidência consultável na interface.
  \end{dxilista}

  \vspace{0.45em}
  {\dxiCorpoP\color{dxinavy!85}Valor esperado: reduzir a navegação entre telas
  e preparar uma síntese verificável sem executar ações operacionais.}
}
\end{dxiconteudo}

\begin{dxiconteudo}{INTEGRAÇÃO · ASSISTENTE DE INVESTIGAÇÃO}
\vspace{-0.40cm}
\dxicol{13.6}{%
  {\dximont\bfseries\fontsize{16}{21}\selectfont\color{lgred}Evolução controlada}

  \vspace{0.3em}
  \begin{dxilista}[dxired]
    \item Fase 1: backend entrega contexto pronto e a IA apenas resume.
    \item Fase 2: seleção restrita entre ferramentas de leitura autorizadas.
    \item Fase 3: piloto com usuários e revisão de utilidade operacional.
    \item Escrita e decisão automática permanecem fora do escopo.
  \end{dxilista}
}%
\dxicol{13.2}{%
  {\dximont\bfseries\fontsize{16}{21}\selectfont\color{hanaroverde}Controles necessários}

  \vspace{0.3em}
  \begin{dxilista}[hanaroverde]
    \item Lista permitida de ferramentas, argumentos e limites de chamadas.
    \item RBAC e organização verificados novamente em cada consulta.
    \item Tempo máximo, telemetria e interrupção segura.
    \item Saída validada antes de chegar à interface.
    \item Histórico de consultas sem conteúdo sensível.
  \end{dxilista}
}
\end{dxiconteudo}

\begin{dxiconteudo}{PLANO PARA VALIDAR AS PROPOSTAS}
\vspace{-0.35cm}
\dxicol{13.8}{%
  {\dximont\bfseries\fontsize{17}{23}\selectfont\color{lgred}Próxima rodada}

  \vspace{0.35em}
  \begin{dxilista}[dxired]
    \item Definir hardware e testar modelos locais mais capazes.
    \item Congelar novos casos sintéticos e gabaritos independentes.
    \item Medir qualidade, evidências, consistência, latência e memória.
    \item Comparar custo e benefício das duas propostas com a baseline.
  \end{dxilista}
}%
\dxicol{13.0}{%
  {\dximont\bfseries\fontsize{17}{23}\selectfont\color{hanaroverde}Decisões solicitadas}

  \vspace{0.35em}
  \begin{dxilista}[hanaroverde]
    \item Qual proposta traz mais valor ao processo do cliente?
    \item Qual perfil participará da validação?
    \item Qual infraestrutura estará disponível para execução local?
    \item Quais limites de tempo e qualidade definem um piloto aceitável?
  \end{dxilista}

  \vspace{0.45em}
  {\dxiCorpoP\color{dxinavy!85}Objetivo do pitch: escolher a prioridade e
  autorizar uma avaliação controlada, sem prometer integração antecipada.}
}
\end{dxiconteudo}

\end{document}
